\section{The SHA-256 Seed Formula}
\label{sec:sha256}

\subsection{Why a Fixed Numerical Seed Is Constitutionally Insufficient}

Early iterations of the \dia{} statute proposed a fixed integer seed---a
specific value published in the statutory text or by the Director of the Census
Bureau before the redistricting cycle begins.
This design is reproducible but exposes a manipulation attack that any
competent opponent will press in litigation.

\begin{description}
\item[The cherry-picking attack.]
If the Director of Census (or any government official with early access to the
census data) can enumerate the partisan consequences of different seeds before
publishing one, the published seed can be selected to advantage one party.
The attack is structurally identical to the ``algorithm selection is a
political choice'' argument established in B.0~\citep{b0paper}: choosing the
seed after knowing its consequences is as political as choosing the algorithm
after knowing its consequences.

\item[The legal consequence.]
Under \rucho{}~\citep{rucho2019}, partisan gerrymandering is a nonjusticiable
federal question.
But the \dia{}'s constitutional claim rests on algorithmic neutrality---the
statute works precisely because no human chose the map.
A fixed seed chosen by an official with knowledge of its partisan consequences
reintroduces the human choice the statute was designed to eliminate.
\end{description}

The SHA-256 formula closes this attack surface completely.

\subsection{The Formula and Its Properties}

\begin{definition}[DIA Seed Protocol v1]
\label{def:diav1}
The seed for the statutory \cs{} run under the \dia{} is
\[
  s_0 = \bigl\lfloor
    \mathrm{SHA\text{-}256}(\mathtt{census\_release\_id}
      \,\|\, \mathtt{"DIA\_SEED\_V1"})
  \bigr\rfloor \bmod 2^{31},
\]
where:
\begin{itemize}
\item $\mathtt{census\_release\_id}$ is the official identifier string for the
      decennial census dataset, published by the Census Bureau under
      42~U.S.C.~\S{}~141 (or successor) at the moment of data release.
      Example: \texttt{"2020-PL94-171-v1.0"}.
\item $\mathtt{"DIA\_SEED\_V1"}$ is the ASCII constant fixed in statute.
      It encodes the version of the seed protocol; a future \dia{} amendment
      adopting a different formula would use \texttt{"DIA\_SEED\_V2"}.
\item $\|$ denotes byte concatenation.
\item SHA-256 is the NIST-standardised cryptographic hash function
      \citep{sha256nist}.
\end{itemize}
\end{definition}

The formula has four properties that jointly satisfy the constitutional
neutrality requirement.

\paragraph{Property~1: Determined by the census, not by any political actor.}
$\mathtt{census\_release\_id}$ is set by the Census Bureau at the moment of
data release and published simultaneously to all parties.
No one---not the party in power, not the Director of Census, not the court
appointing a special master---can know $s_0$ until after the census data is
released.
The seed is discovered, not chosen.

\paragraph{Property~2: The protocol constant is version-locked in statute.}
$\mathtt{"DIA\_SEED\_V1"}$ is a literal string in the statutory text.
Neither Congress nor the Executive can change it without a formal statutory
amendment, which is subject to the full legislative process.
This prevents administrative modification of the seed formula.

\paragraph{Property~3: The seed changes automatically with the census.}
SHA-256 is a cryptographic hash: a one-character change in the input produces
a completely different 256-bit output.
The 2020 census release identifier differs from the 2010 identifier; therefore
$s_0^{2020} \ne s_0^{2010}$.
Each census cycle produces an independent seed, preventing any carry-over
advantage from one cycle to the next.

\paragraph{Property~4: No one can predict the seed before the census is published.}
SHA-256 is computationally preimage-resistant: given $s_0$, no efficient
algorithm can recover $\mathtt{census\_release\_id}$, and given a partial
census, no efficient algorithm can predict the final
$\mathtt{census\_release\_id}$ well enough to precompute $s_0$.
The partisan consequences of $s_0$ are unknowable until after the census
data---and therefore the redistricting problem itself---is fully specified.

\subsection{Why the Version String Is Necessary}

Without the version string \texttt{"DIA\_SEED\_V1"}, the formula would be
$s_0 = \lfloor \mathrm{SHA\text{-}256}(\mathtt{census\_release\_id}) \rfloor
\bmod 2^{31}$.
This has an undesirable property: any other application that hashes the census
release identifier (for any purpose) produces the same value.
If the Census Bureau or any government agency independently publishes a SHA-256
hash of the release identifier as a data-integrity check, an adversary could
use that published hash to precompute $s_0$ before the redistricting cycle
begins.

The version string \texttt{"DIA\_SEED\_V1"} domain-separates the
redistricting seed from all other SHA-256 uses of the same input.
It is standard cryptographic hygiene (domain separation)~\citep{bernstein2019}.

\subsection{The Seed Is Not Secret}

A common misunderstanding is that the SHA-256 formula is designed to hide the
seed from adversaries.
It is not.
$s_0$ is public information: the moment the Census Bureau releases the census
data, any party can compute $s_0$ in milliseconds.

The protection the formula provides is \emph{temporal}, not \emph{epistemic}.
The seed cannot be known before the census is published.
Once the census is published, $s_0$ is universally known, the redistricting
run is public, and the resulting map is auditable by any party.
This is the correct design for a democratic institution: the process is opaque
before the data exist, and fully transparent afterward.

\subsection{Litigation Posture}

The SHA-256 formula allows the \dia{} to make the following affirmative claim
in any court challenge:

\begin{statute}
\textbf{Certificate of Seed Neutrality.}

The redistricting seed used in this proceeding is
$s_0 = \lfloor \mathrm{SHA\text{-}256}(\texttt{census\_release\_id} \,\|\,
\texttt{"DIA\_SEED\_V1"}) \rfloor \bmod 2^{31}$.

\medskip

This value was not known to any party until the Census Bureau published the
dataset identifier $\mathtt{census\_release\_id}$ on the date of record.
No human selected this seed; it was computed mechanically from public data.
The partisan consequences of this seed were not knowable before that date.

\medskip

The seed is reproducible: any party, any court, any special master may
independently compute $s_0$ from the public census release identifier and
verify that the \cs{} run with this seed and threshold $T = 600$ produced the
map now before the court.
\end{statute}

No fixed-integer seed can make this certificate.
The SHA-256 formula can.
